\section*{Bab \ref{ch:g7:sediments-and-fossils} --- Sedimen dan Fosil}

\begin{solution}{exo:g7:sediments-and-fossils:1}
(a) \omterm{def:g7:sediments-and-fossils:fossil}{cetakan}; (b) \omterm{def:g7:sediments-and-fossils:fossil}{fosil jejak}; (c) \omterm{def:g7:sediments-and-fossils:sediment}{sedimen}; (d) \omterm{def:g7:sediments-and-fossils:fossil}{tuangan}; (e) \omterm{def:g7:sediments-and-fossils:layer}{stratum};
(f) \omterm{def:g7:sediments-and-fossils:fossil}{fosil}.
\end{solution}

\begin{solution}{exo:g7:sediments-and-fossils:2}
(b) hewan itu mati; (d) \omterm{def:g3:stones-pebbles-sand:sand}{lumpur} cepat menutupi cangkang itu; (c) lapisan-lapisan
baru menumpuk dan cangkang itu berubah menjadi batu; (a) \omterm{def:g5:water-shapes-the-land:erosion}{erosi} mengikis
lapisan-lapisan di atasnya.
\end{solution}

\begin{solution}{exo:g7:sediments-and-fossils:3}
(a) \omterm{def:g2:water-on-earth:ocean}{Laut} yang hangat, dangkal, dan jernih. (b) Air yang dangkal, di tepi
\omterm{def:g1:landscapes-around-us:water}{danau}, \omterm{def:g1:landscapes-around-us:water}{sungai}, atau \omterm{def:g2:water-on-earth:ocean}{laut}. (c) Hutan rawa atau paya. (d) Daratan yang
ditumbuhi tumbuhan (dan tempat capung terbang), di dekat air tempat \omterm{def:g3:stones-pebbles-sand:sand}{lumpur}
itu mengendap.
\end{solution}

\begin{solution}{exo:g7:sediments-and-fossils:4}
(a) Salah: sebagian besar makhluk hidup tidak meninggalkan \omterm{def:g7:sediments-and-fossils:fossil}{fosil}; makhluk
itu membusuk atau larut sebelum terkubur. (b) Benar: jejak kaki adalah
\omterm{def:g7:sediments-and-fossils:fossil}{fosil jejak}. (c) Salah: butir-butir yang paling besar dan paling berat
mengendap lebih dahulu, \omterm{def:g3:stones-pebbles-sand:sand}{lumpur} halus paling akhir. (d) Benar.
\end{solution}

\begin{solution}{exo:g7:sediments-and-fossils:5}
(a) Batu pasir \qty{120}{m}, batu kapur \qty{180}{m}, serpih \qty{100}{m}.
(b) $120 + 180 + 100 = \qty{400}{m}$. (c) $180 \div 400 = 0.45$:
\textbf{45~\%}. (d) Puncaknya telah terkikis oleh \omterm{def:g5:water-shapes-the-land:erosion}{erosi}: yang tersisa bukan
seluruh lapisan, dan banyaknya yang tersisa berbeda-beda dari tempat ke
tempat.
\end{solution}

\begin{solution}{exo:g7:sediments-and-fossils:6}
(a) Dari bawah ke atas: \omterm{def:g3:stones-pebbles-sand:pebble}{kerikil}, lalu \omterm{def:g3:stones-pebbles-sand:sand}{pasir}, lalu \omterm{def:g3:stones-pebbles-sand:sand}{lumpur} di atasnya, di
bawah air yang jernih. (b) Butir-butir \omterm{def:g3:stones-pebbles-sand:sand}{lumpur} yang sangat kecil tenggelam
dengan sangat lambat: butir-butir itu masih berjatuhan sesudah satu jam.
(c) Seperangkat lapisan kedua di atas yang pertama: \omterm{def:g3:stones-pebbles-sand:sand}{pasir}, lalu \omterm{def:g3:stones-pebbles-sand:sand}{lumpur}
lagi. (d) \omterm{def:g3:stones-pebbles-sand:pebble}{Kerikil} di bagian paling bawah.
\end{solution}

\begin{solution}{exo:g7:sediments-and-fossils:7}
Paling berpeluang: (b), cangkang yang keras, cepat terkubur di bawah \omterm{def:g3:stones-pebbles-sand:sand}{lumpur}
halus di air yang tenang. Lalu (d): cangkangnya keras dan terkubur, tetapi
\omterm{def:g5:water-shapes-the-land:flood}{banjir} dapat juga memindahkan dan memecahkannya. Paling kecil peluangnya:
(a), tubuh yang lunak tanpa bagian keras, di atas batu-batu tempat tidak ada
\omterm{def:g3:stones-pebbles-sand:sand}{lumpur} yang mengendap untuk menutupinya. (c) tidak jauh lebih baik: tulang
itu tetap tidak tertutup lalu terurai dan terkikis.
\end{solution}

\begin{solution}{exo:g7:sediments-and-fossils:8}
(a) Sebuah rongga dengan bentuk cangkang tercetak di dindingnya: sebuah
\omterm{def:g7:sediments-and-fossils:fossil}{cetakan}. (b) Sebuah \omterm{def:g7:sediments-and-fossils:fossil}{tuangan}, salinan cangkang dari batu. (c) Tidak:
cangkangnya sudah larut; \omterm{def:g7:sediments-and-fossils:fossil}{tuangan} itu terbuat dari \omterm{def:g6:rocks-and-minerals:mineral}{mineral} baru. (d) Air
mudah lewat di antara butir-butir \omterm{def:g3:stones-pebbles-sand:sand}{pasir} yang besar, sehingga air dapat
mencapai dan melarutkan cangkang itu; di dalam \omterm{def:g3:stones-pebbles-sand:sand}{lumpur} halus, air jauh
lebih sulit lewat.
\end{solution}

\begin{solution}{exo:g7:sediments-and-fossils:9}
(a) $1.6 \times 1000 = \qty{1600}{m}$. (b) $1600 \div 80 = 20$ lapisan.
(c) $160 \div 1600 = 0.1$: \textbf{10~\%}.
\end{solution}

\begin{solution}{exo:g7:sediments-and-fossils:10}
(a) Tidak: kaki hanya meninggalkan jejak di tanah yang lunak, \omterm{def:g3:stones-pebbles-sand:sand}{pasir} basah,
atau \omterm{def:g3:stones-pebbles-sand:sand}{lumpur}; tanah itu baru menjadi \omterm{def:g3:stones-pebbles-sand:rock}{batuan} kemudian. (b) Ombak atau \omterm{def:g5:water-shapes-the-land:flood}{banjir}
berikutnya menyapunya. (c) \omterm{def:g7:sediments-and-fossils:sediment}{Sedimen} baru harus cepat menutupi dan mengisinya,
sebelum jejak itu terhapus. (d) \omterm{def:g7:sediments-and-fossils:fossil}{Fosil jejak}: \omterm{def:g7:sediments-and-fossils:fossil}{fosil} itu menunjukkan apa yang
dilakukan hewan itu, bukan bagian dari tubuhnya.
\end{solution}

\begin{solution}{exo:g7:sediments-and-fossils:11}
(a) $20 + 5 = \qty{25}{mm}$. (b) $1000 \div 25 = 40$ pasang: \omterm{def:g6:inside-the-earth:layers}{inti} bor itu
mencatat \textbf{40 \omterm{def:g5:water-shapes-the-land:flood}{banjir}}. (c) \omterm{def:g3:stones-pebbles-sand:sand}{Pasir} itu berat: \omterm{def:g3:stones-pebbles-sand:sand}{pasir} mengendap tidak
lama sesudah \omterm{def:g1:landscapes-around-us:water}{sungai} memasuki \omterm{def:g1:landscapes-around-us:water}{danau}, di dekat muara; hanya sedikit yang
terbawa jauh. (d) Jauh di tengah, hanya \omterm{def:g3:stones-pebbles-sand:sand}{lumpur} halus yang masih ada di
dalam air: \omterm{def:g3:stones-pebbles-sand:sand}{lumpur} itu mengendap di sana perlahan-lahan, ketika airnya
tenang.
\end{solution}

\begin{solution}{exo:g7:sediments-and-fossils:12}
(a) Cangkang itu keras: cangkang bertahan cukup lama untuk terkubur, dan
banyak hewan bercangkang hidup di dasar \omterm{def:g2:water-on-earth:ocean}{laut}, tempat \omterm{def:g3:stones-pebbles-sand:sand}{lumpur} mengendap di
atasnya. (b) Tubuh yang lunak membusuk sebelum sempat terkubur. (c) Tidak:
hewan bertubuh lunak hampir tidak pernah meninggalkan \omterm{def:g7:sediments-and-fossils:fossil}{fosil}, sehingga
daftar dari \omterm{def:g7:sediments-and-fossils:fossil}{fosil} tidak lengkap; daftar itu terutama menunjukkan hewan yang
bercangkang atau bertulang.
\end{solution}

\begin{solution}{pb:g7:sediments-and-fossils:1}
\textbf{1.} Batu pasir, batu kapur, serpih, kapur (yang pertama di bagian
bawah kolom).

\textbf{2.} \qty{120}{m}, \qty{180}{m}, dan \qty{100}{m}; totalnya
\qty{400}{m}.

\textbf{3.} $180 \div 400 = 0.45$: \textbf{45~\%}.

\textbf{4.} Puncaknya telah terkikis: bagian atas lapisan itu sudah hilang,
dan tidak ada yang dapat mengukur apa yang sudah tidak ada lagi.

\textbf{5.} Di air yang dangkal di tepi \omterm{def:g1:landscapes-around-us:water}{danau}, \omterm{def:g1:landscapes-around-us:water}{sungai}, atau \omterm{def:g2:water-on-earth:ocean}{laut}: gelembur
gelombang terbentuk di \omterm{def:g3:stones-pebbles-sand:sand}{pasir} di air yang dangkal, dan jejak kaki
ditinggalkan di \omterm{def:g3:stones-pebbles-sand:sand}{pasir} atau \omterm{def:g3:stones-pebbles-sand:sand}{lumpur} yang lunak dan basah, sering di pantai.

\textbf{6.} Di \omterm{def:g2:water-on-earth:ocean}{laut} yang hangat, dangkal, dan jernih: karangnya.

\textbf{7.} Tenang: \omterm{def:g3:stones-pebbles-sand:sand}{lumpur} halus hanya mengendap di tempat airnya hampir diam.

\textbf{8.} Mula-mula tempat itu adalah pantai, dengan air yang dangkal dan
\omterm{def:g3:stones-pebbles-sand:sand}{pasir} yang lunak tempat hewan-hewan berjalan. Lalu \omterm{def:g2:water-on-earth:ocean}{laut} menutupinya: \omterm{def:g2:water-on-earth:ocean}{laut}
yang hangat dan dangkal tempat karang tumbuh dan cangkang menumpuk.
Kemudian airnya menjadi lebih tenang, mungkin lebih dalam, dan hanya \omterm{def:g3:stones-pebbles-sand:sand}{lumpur}
halus yang mengendap, bersama cangkang amonit, hewan yang berkerabat dengan
cumi-cumi dan gurita.

\textbf{9.} $180 \times 100 = 18\,000$~cm.

\textbf{10.} $120~\text{m} = 12\,000$~cm, dan $12\,000 \times 20 =
240\,000$~tahun.

\textbf{11.} $1~\text{m} = 100$~cm, dan $100 \times 100 = 10\,000$~tahun.

\textbf{12.} $180 \times 10\,000 = 1\,800\,000$: kira-kira
\textbf{1\,800\,000 tahun}, hampir dua juta tahun. Itu
$1\,800\,000 \div 240\,000 = 7.5$ kali lebih lama daripada batu pasir.
\end{solution}
